%% 第31回編集委員会（物理第5巻）：事務局（前付・巻末の組）
%% 物理シリーズ　第5巻　原子　本文（ph4-body.tex と同じ形）
%% 章は \bkinputchap{p5-NN} で読む（原-0 が p5-00、原-16 が p5-16。まだない章は、組むときに飛ばして「BKPAGE missing」を出す）
%% 部は \phPartPages{番号}{問い}{デルタのせりふ}{この部を読むと…}：部扉（左）・この部で学ぶこと（右。ph5-part<番号>.tex）・柱とつめ
%%   部の名前は phys/series_phys.py の第5巻（ph5-meta.tex の \phdefpart）。第4部は原-11〜原-16。
%%   部扉の3つの問いは、その部の章の「なぜ？」（もくじの問いと、章の中の「まだわかっていないこと」）から選んだ。
%%   裏表紙の5つ（原-4・3・9・12・13：gen_phys.py の WHY_PICK）とは重ねない。第1部の3つめは原-2 の「まだわかっていないこと」、
%%   第3部の3つめは原-10 の「まだわかっていないこと」。第4部は原-11・原-14・原-15 のもくじの問い。
%%   デルタのせりふと「この部を読むと」は、各章のねらい（goal）と節の見出しから書いた。
%%   第4部のせりふは、N-6・N-7（「安全か危険か」の二択にせず量で書く。賛成・反対の立場をとらない）に合わせた。
% 組版のメモリ対策（全巻にかかる修正事項 No.58。第9巻と同じ）：章の前に Lua のごみ集めをする
\providecommand\bkinputchap[1]{\directlua{collectgarbage("collect") texio.write_nl("log","BKMEM #1 "..math.floor(collectgarbage("count")/1024).." MB")}\IfFileExists{#1.tex}{\input{#1}}{\clearpage\typeout{BKPAGE missing #1}}}
\bkinputchap{p5-00}

%% ======================== 第1部 飛び出す：光と電子の二重性 ========================
%% 部の名前の終わりの字（「性」「口」「体」）が、うすい大きな数字にかかる（check_overlap の候補）ので、第1〜3部は数字を右へ寄せる
%% （第4巻の第2部と同じやり方。共通部品は変えない）
\makeatletter\def\bk@partnumpos{164,212}\makeatother
\phPartPages{1}%
  {172/{原子の中は、ほとんどからっぽだって本当？},
   160/{熱した鉄は、なぜ温度で光の色が変わる？},
   148/{光が当たってから、電子が飛び出すまでの時間は？}}%
  {原子をのぞくと、まん中に小さくて重い原子核、周りに軽い電子がいるよ。光はつぶとしても飛びこみ、電子は波としても広がる。原子が決まった色の光だけを出すのは、エネルギーがとびとびの段になっているからなんだ。}%
  {電子の比電荷と電気素量をはかる実験と、α線の散乱から原子の形を読み、光電効果とコンプトン効果で光子のエネルギー $E=hf$ と運動量を求め、電子の波長から電子顕微鏡と不確定性関係を説明し、ボーアの模型で水素のスペクトル・X線・レーザーのしくみを計算できるようになります。この部の実験の事実は、第2部から先の土台です。}
\bkinputchap{p5-01}
\bkinputchap{p5-02}
\bkinputchap{p5-03}
\bkinputchap{p5-04}

%% ======================== 第2部 すりぬける：量子力学の入口 ========================
\phPartPages{2}%
  {172/{電子は、原子の中のどこにいる？},
   160/{USBメモリは、どうやって電子を閉じこめて記録する？},
   148/{ネオンは、なぜほかの物質とほとんど結びつかない？}}%
  {電子の波は、弦と同じように、決まった形の定在波だけが残るよ。壁があっても、波は少しだけしみ出して、向こう側へすりぬける。原子の中の電子の「席」の決まりから、周期表の形まで出てくるんだ。}%
  {箱に閉じこめた電子のエネルギーを定在波から求め、シュレディンガー方程式を立てて解き、$|\psi|^2$ を見つかる確率として読み、トンネル効果でフラッシュメモリと走査トンネル顕微鏡のしくみを説明し、調和振動子で分子のゆれを、量子数で電子の席を数えて周期表と化学結合を説明できるようになります。}
\bkinputchap{p5-05}
\bkinputchap{p5-06}
\bkinputchap{p5-07}

%% ======================== 第3部 切りかえる：固体と半導体 ========================
\phPartPages{3}%
  {172/{金属は電気を通すのに、ガラスはなぜ通さない？},
   160/{スマホのチップは、どうやって計算している？},
   148/{トランジスタは、どこまで小さくできる？}}%
  {原子がたくさん並ぶと、電子のエネルギーの段は、はばのある帯になるよ。帯のすき間の大きさで、金属・半導体・絶縁体が分かれる。半導体を組み合わせると、電流を一方通行にしたり、オンとオフを切りかえたりできるんだ。}%
  {X線の回折で結晶をはかり、フェルミ分布とエネルギーの帯で金属・半導体・絶縁体の違いを説明し、n形・p形の半導体と pn 接合でダイオード・LED・太陽電池のしくみを導き、トランジスタをスイッチと増幅器として使い、CMOS の論理ゲートとオペアンプの回路を計算できるようになります。}
\bkinputchap{p5-08}
\bkinputchap{p5-09}
\bkinputchap{p5-10}

%% ======================== 第4部 分かれる：原子核と放射線 ========================
\makeatletter\def\bk@partnumpos{152,212}\makeatother% 第4部はもとの位置（bk-physbody.sty の値）にもどす
\phPartPages{4}%
  {172/{太陽は、何をエネルギーのもとにして光っている？},
   160/{目に見えない放射線を、どうやってはかる？},
   148/{原子炉は、どうやって反応を一定に保つ？}}%
  {原子核の中では、質量の一部がエネルギーに変わるよ。こわれる原子核からは放射線が出る。放射線の影響は「あるか、ないか」ではなく、受けた量で決まるんだ。量をはかり、比べて考えよう。}%
  {相対論から $E=mc^2$ を導いて結合エネルギーと核反応のエネルギーを計算し、半減期と崩壊の法則で放射能と年代を求め、放射線が物を通るときの減り方と遮へいを計算し、Bq・Gy・Sv で線量をはかって比べ、増倍率と遅発中性子で原子炉の制御と崩壊熱を説明し、原子の力とのつきあい方を、物理の量と公的な資料の数で考えられるようになります。}
\bkinputchap{p5-11}
\bkinputchap{p5-12}
\bkinputchap{p5-13}
\bkinputchap{p5-14}
\bkinputchap{p5-15}
\bkinputchap{p5-16}

%% ======================== 巻末 ========================
%% 第31回編集委員会（物理第5巻）：事務局（前付・巻末の組）
%% 物理シリーズ　第5巻　原子　巻末（こたえ・記号カード・核種のデータ・物性値・定数・なぜ？の道すじ・チェック表・つぎの巻へ・さくいん・奥付）
%% 第4巻の ph4-back.tex と同じ部品・同じ形。見出し（\bkbacktitle）の名前は、もくじ（ph5-toc.tex）の巻末の項目の名前と同じにする
%%   （tools/check_pages.py p5 が照らし合わせる。表の名前は phys/series_phys.py の第5巻の tables「核種のデータ・物性値・定数」。
%%   構成案の巻末「核種のデータと減弱係数」と、計画書 C-3 の名前）。
%% 第1〜4巻の確認で決めたこと（U-1〜U-13、Q6、X-1〜X-8、Y-1〜Y-12、G-1〜G-15）を、はじめから当てはめた：
%%   ・記号カード：数学の記号の行は、この巻の章の \bkkigocard（自動でうしろに集まる）と同じ記号・同じ意味を置かない（U-6・X-2）。
%%   ・紙面に委員会の回を書かない（U-7）。回はこのコメントにだけ書く。
%%   ・さくいんの段は \raggedcolumns。奥付の「先に読む本」は表4（bk-phys.sty）と同じ文（Q6-22）。
%%   ・「分野で記号がぶつかるもの」は、X-8 のやり方（同じ記号を2行にしない、章で確かめた意味だけを書く）で、11月11日の17章の原稿の
%%     頭のコメント（記号の欄）・記号カード・本文で確かめて作った。Z5-3 の μ・κ・A・B・λ と、原-13 の中性子の減速の M
%%     （docs/作業メモ/物理第5巻_組どうしの確かめ.md の「ほかの組・共通部品への依頼」）、計画書 E の4 の E・σ・k・ρ・φ・D・N・n・β・γ の行を入れた。
%%   第32回（修5-6）：F6-32 と H-6 で、章で分けた記号（原-6 の A、原-7 の M_モル・n_原子・n_電子、原-8 の r_F、原-10 の帰還率 β、原-13 の I_ex・n_e、
%%     原-14 の計数率 m と体の質量 M、原-15 の炉の周期 T・α_T・γ_I・γ_X・I・利得 Q・段の数 m、原-15・16 の Δθ）と θ の行を足した。
%%     ほかの組（修5-1〜5）の直しと同時なので、まとめ役が直したあとの原稿と照らし合わせる。
%%   第32回（巻を組む組、11月12日）：直したあとの17章の原稿の数式の中の記号を章ごとに数え、行に書いた章と、記号が出る章を比べて行を合わせた。
%%     A（原-6 ⑥の係数は c_1・c_2 になり、⑪の行き来の強さだけが残る。原-5 ④の波動関数の大きさ）、M・m（原-8 も M_モル。原-4 ⑤の陽子の質量・
%%     原-11 ④の原子の質量）、ρ（原-0・7・8・11 の密度。原-8 は抵抗率を記号にしない）、N・n（原-1 の数密度は n、個数の N、原-2・6 の量子数、
%%     原-8 ②の反射の次数、原-14 の本当の計数率 n）、R（原-1・3・13 の半径・飛程）、d・Δx の行を新しく作った（原-5 ⑫の点の間かくは d になり、
%%     不確かさ Δx とぶつからない。原-3 ⑤のしまの間かく Δx は波-13 の式を引いたもの）。ほか τ・α・β・γ・η・ε・Q・S・a・L・I・g・H・k・ν・σ に足した。
%%     照らし合わせは checks/物理第5巻/巻末/symtab_ph5.py（出力 symtab_ph5_out.txt。残りは同じ意味か、表計算のセル・図の札・単位）。
%%     行が長くなったので、表の見出しの前の \Needspace を 18→12（H-1）、列の幅を 15・34・28→14・31・25 mm、rowsep を 1.4→1.0pt にした
%%     （約束の最後の2項目だけが次のページに残っていたのが消え、記号カードは 16→14 ページ）。
%% 表の値は、各章のファイルの頭のコメント（出典の欄）と、docs/作業メモ/物理第5巻_組どうしの確かめ.md の「巻末の表に入れる値」
%%   「出典の組がしたこと」から集めた（出典つきの値だけ）。計算は checks/物理第5巻/巻末/calc_ph5_back.py（出力 calc_ph5_back_out.txt）。
%%   ・定数：CODATA 2022（NIST allascii.txt。原-1〜原-11 の組が 11月11日に curl で読んだ値）。導いた値（ħc・ke²・hc・k_BT）は計算し直した。
%%   ・原子の質量：NIST（AME2020。checks/物理第5巻/原-11/nist_masses_20261111.json）。1核子あたりの結合エネルギーと Q 値は計算し直した
%%     （Q 値は ENSDF の Ra-226 4.8707・Po-210 5.4075・U-238 4.2699・Po-212 8.954 MeV と合う）。
%%   ・半減期：ENSDF の評価値（IAEA LiveChart の API。出典の組が 11月11日に確かめた）と環境省の基礎資料（令和7年度版）。
%%   ・〔要確認〕は、章の頭のコメントで〔要確認〕のままの値：ケイ素・ゲルマニウム・GaAs・GaN の物性値、銅のデバイ温度、Nb-Ti の臨界温度、
%%     鉄の 2.2 μ_B、STM の仕事関数 4.5 eV、遅発中性子の割合 β＝0.0065 の ν_d、崩壊熱の式の係数 0.066、Cs-137 の 1 cm 線量当量率定数 0.0927。
%% 「つぎの巻へ」：第5巻はシリーズの最後の巻なので、見出しの名前（もくじの項目）は前の巻と同じ「つぎの巻へ」のまま、中身を
%%   「全5巻のふりかえり」と「この先の学び（工学の専門の科目）」にした。表は、各章の「工学の駅」（\phstation）の科目を集めた
%%   ph5-nextuse.tex（checks/物理第5巻/巻末/mknext_p5.py が作る。tools/mknext.py p5 では、後ろの巻がないので行が0になるため）。
%% 表の書式（見出しの行の色・横の線・行の間）は、第4巻と同じく、各表の colspec のあとにじかに書いた。
\makeatletter
\def\phbkhead#1{\par\addvspace{4mm}\Needspace{6\baselineskip}\noindent{\sffamily\bfseries\large\color{bkDeep}#1}\par\nobreak\vspace{2mm}}
\newcommand\phyk{〔要確認〕}
%% ---------------- こたえ ----------------
\bkAnswerPages
\par\vspace{4mm}
%% 箱はページをまたがせない（またぐと、次のページに箱の切れはしだけが残る）
\begin{tcolorbox}[bk base, unbreakable, colframe=bkTheme!60, colback=bkTheme!5, boxrule=0.8pt, arc=2mm, left=4mm, right=4mm, top=2mm, bottom=2mm]
{\sffamily\bfseries\color{bkDeep}答え合わせのしかた}\par\vspace{1mm}
\begin{enumerate}[label={\bkcircled{\arabic*}}, leftmargin=2.4\zw, itemsep=0pt, topsep=0pt]
\item まず自分の答えを全部書いてから、ここを見よう。
\item 単位と桁を確かめよう。J と eV（$\qty{1}{eV}=\qty{1.60e-19}{J}$）、nm・pm・fm、cm$^{-3}$ と m$^{-3}$、Bq・Gy・Sv を取り違えていないかも。1年は365.25日で計算しています。
\item 丸め方で最後の桁が1つずれても、間違いではありません。$h=6.63\times10^{-34}$ と $6.626\times10^{-34}$、$hc=\qty{1240}{eV.nm}$ と $1239.8$ のどちらで計算しても、2けたは同じになります。
\item 違ったら、同じ形の例題にもどって、どの段で違ったかを探そう。$\lambda$ が波長か崩壊定数か、$A$ が質量数か放射能か、$E$ がエネルギーか実効線量かも確かめよう。
\end{enumerate}
\end{tcolorbox}

%% ---------------- 記号カード ----------------
%% 数学シリーズの記号で、この巻でよく使うもの（この巻の章の \bkkigocard は自動でうしろに集まる）。
%% 章のカードにある記号（11月11日の原稿で数えた。原-1：e/m、Z、b、d、A（第32回 H-7）　原-2：h、W、K0、V0　原-3：ħ、Δx、ψ、v_g　原-4：n、R∞、a0、A・B
%%   原-5：i、⟨x⟩、m*　原-6：κ、𝒯・ℛ、μ、ν̃　原-7：l、μ_B、Z_eff　原-8：E_F、E_g、θ_D、κ　原-9：N_d、𝓔、μ_n　原-10：g_m、h_FE、V_th、α
%%   原-11：γ、B、u、Q　原-12：λ、T_1/2、A　原-13：μ、x_1/2、σ、Σ、φ、S、ξ　原-14：D、H_T、E、w_R、w_T、τ、Γ　原-15：k、ρ、ν、β、B²）は、
%%   同じ意味ではここに置かない（U-6）。虚数単位 i は原-5 のカードなので、数学の行は回路の章（原-9・10）の j だけにした。
%%   ψ* の共役は 6-8（数学の本の上線）。固有値の λ は 7-8。
\bkKigoPages{%
\bkkigorow{$a\times10^{n}$}{エー・かける・じゅうのエヌじょう}{とても大きい数・小さい数の書き方。$\qty{1}{fm}=10^{-15}$\,m、$N_{\mathrm A}=6.02\times10^{23}$}{2-7}%
\bkkigorow{$\sqrt{x}$}{ルート}{平方根。ローレンツ因子 $1/\sqrt{1-v^2/c^2}$、数えた数のばらつき $\sqrt N$}{2-8}%
\bkkigorow{$e^x$、\ $\ln x$}{イー、エル・エヌ・エックス}{ネイピア数（2.718…）と自然対数。崩壊 $e^{-\lambda t}$、減弱 $e^{-\mu x}$、すりぬけ $e^{-2\kappa a}$。電気素量の $e$ とは別}{3-11}%
\bkkigorow{$\log_{10}x$}{ログ}{常用対数。片対数のグラフ、ダイオードの「約 60 mV で電流10倍」}{3-12}%
\bkkigorow{$\sigma$（ばらつき）}{シグマ}{標準偏差。数えた数 $N$ のばらつきは約 $\sqrt N$（原-12 ⑨）。断面積の $\sigma$（原-13）とは別}{3-14}%
\bkkigorow{$\approx$}{ニアリーイコール}{ほぼ等しい（近似）}{5-8}%
\bkkigorow{$\displaystyle\int$、\ $\displaystyle\sum$}{インテグラル、シグマ}{細かく分けてたす・全部たす。$\displaystyle\int\lvert\psi\rvert^2\dd x=1$。$\sum$ は巨視的断面積 $\Sigma$（原-13）とは別}{5-10}%
\bkkigorow{$j$}{ジェイ}{虚数単位（2乗すると $-1$）。回路の章（原-9・10）で使う。原-5〜8 の $i$（原-5）と同じ数}{6-8}%
\bkkigorow{$\bar z$、\ $\psi^*$}{ゼット・バー、プサイ・スター}{共役な複素数（虚部の符号を変えた数）。数学の本の上線を、量子力学では $^*$ で書く（原-5）}{6-8}%
\bkkigorow{$e^{j\theta}$}{イー・ジェイ・シータ}{単位円の上の角 $\theta$ の点（オイラーの公式）。量子力学では $e^{i\theta}$、時間の部分 $e^{-iEt/\hbar}$（原-5）}{6-10}%
\bkkigorow{$\sinh x$}{ハイパボリック・サイン}{双曲線関数。壁をすりぬける正確な透過率（原-6）}{6-11}%
\bkkigorow{$\mathbf A\vec v=\lambda\vec v$}{固有値}{固有値と固有ベクトル。とびとびのエネルギーはシュレディンガー方程式の固有値（原-5 ⑦）。波長・崩壊定数の $\lambda$ とは別}{7-8}%
\bkkigorow{$\tfrac{\partial f}{\partial x}$}{ラウンドディー}{偏微分。時間をふくむシュレディンガー方程式（原-5 ⑩）}{8-2}%
\bkkigorow{$\nabla^2$}{ラプラシアン}{3次元のシュレディンガー方程式（原-7）、1群の拡散方程式（原-15）}{8-12}%
}

\Needspace{12\baselineskip}\phbkhead{分野で記号がぶつかるもの}
{\small 原子の物理は、量子・固体・回路・原子核・放射線と分野が広く、同じ文字が違う量を表すことがあります。この本では、同じ章の中では1つの文字を1つの意味にだけ使い、ぶつかるときは文字や添字・書体を変えました。どの量かは、章（部）と単位で見分けよう。右から2つめの列は、第1巻『力学』〜第4巻『電磁気学』とほかの巻での意味です。\par}\vspace{2mm}
{\footnotesize
\begin{longtblr}{colspec={Q[c,m,wd=14mm] X[l,m] Q[l,m,wd=31mm] Q[l,m,wd=25mm]}, width=\linewidth, rowhead=1, row{1}={bg=bkVol!14, font=\sffamily\bfseries\footnotesize}, hlines={0.4pt,bkLine}, rowsep=1.0pt, colsep=3.5pt}
記号 & この巻での意味 & 第1〜4巻・ほかの巻での意味 & 見分け方 \\
$h$ & プランク定数〔J·s〕（全章。$\hbar=h/(2\pi)$ は原-3 から）、直流電流増幅率 $h_{\mathrm{FE}}$（原-10） & 高さ（力）、比エンタルピー・熱伝達率（熱） & 単位：J·s。$h_{\mathrm{FE}}$ は単位なし \\
$f$、$\nu$ & 光の振動数 $f$〔Hz〕（$E=hf$。高校の教科書の $h\nu$）、電波・回路の周波数 $f$（原-7・10）、4因子の熱中性子利用率 $f$（原-15）。$\nu$ は1回の分裂で出る中性子の数（原-15）、分光の波数は $\tilde\nu$〔cm$^{-1}$〕（原-6）。ニュートリノ $\nu$・$\bar\nu$（原-11・12）は粒子の記号 & 振動数（波）、周波数（電） & 単位：Hz、cm$^{-1}$。$\nu$ は単位なし \\
$\lambda$ & 波長〔m〕（第1部・原-8・9・11）、崩壊定数〔1/s〕（第4部の原-12〜15） & 波長（波・電-13・15）、線電荷密度（電-1）、管摩擦係数（力-14）、固有値（7-8） & 部と単位：m、1/s \\
$E$、$\mathcal E$ & エネルギー〔J・eV〕（全章）、電場 $E$（原-1 だけ）、直流の電源の電圧 $E$（原-10）、実効線量 $E$〔Sv〕（原-14。エネルギーは $E_\gamma$ と添字）。帯の端 $E_{\mathrm c}$・$E_{\mathrm v}$、フェルミエネルギー $E_{\mathrm F}$、禁制帯幅 $E_{\mathrm g}$（原-8・9）。半導体の中の電場は書体を変えて $\mathcal E$（原-9） & 電場（電）、起電力 $\mathcal E$（電-4・5・8）、ヤング率（力-8）、照度（波-12） & 単位：J・eV、V/m、V、Sv \\
$k$、$K$ & 波数 $k$〔1/m〕（原-3・5・6・8）、クーロンの定数 $k$（原-1・4・7・11）、増倍率 $k$（原-15）、ばね定数 $k_{\mathrm s}$（原-6）、MOSFET の係数 $k'$（原-10）、さいころをふる回数 $k$（原-12）。運動エネルギー $K$（原-1・2・11）、結晶の中の波の番号 $K$（原-8）。ボルツマン定数は必ず $k_{\mathrm B}$ & ばね定数（力-2）、熱伝導率（熱-11）、波数（波-5） & 単位：1/m、N·m$^2$/C$^2$。増倍率は単位なし \\
$\mu$ & 換算質量〔kg〕（原-6・原-7 ⑪）、減弱係数〔1/cm〕（原-13・14。質量減弱係数は $\mu/\rho$）、ポアソン分布の平均（原-12）。移動度 $\mu_{\mathrm n}$・$\mu_{\mathrm p}$（原-9・10）、ボーア磁子 $\mu_{\mathrm B}$（原-7）。単位の μ（マイクロ）は立てた字 & 透磁率 $\mu_0$（電）、摩擦係数・粘度（力）、線密度（波-5）、化学ポテンシャル（熱-10・15。原-8 では $E_{\mathrm F}$） & 添字と単位：kg、1/cm、cm$^2$/(V·s) \\
$\kappa$ & 減衰の定数〔1/m〕（原-6）、熱伝導率〔W/(m·K)〕（原-8） & 熱伝導率は $k$（熱-11） & 単位：1/m、W/(m·K) \\
$A$ & 質量数（原-1・11。原-12〜14 では核種の記号の左上の数としてだけ書く）、放射能〔Bq〕（原-12〜14）、アインシュタインの係数 $A$（原-4 ⑫）、波動関数の大きさ $A$（原-5 ④）、トンネル分裂の行き来の強さ $A$〔eV〕（原-6 ⑪。⑥の井戸の係数は $c_1$・$c_2$）、オペアンプの増幅度 $A$（原-10） & 面積（電-14〜16）、振幅（波）、四端子定数（電-13） & 章と単位：Bq。質量数は単位なし \\
$B$ & 磁場（磁束密度）$\vec B$〔T〕（原-1・7）、原子核の結合エネルギー〔MeV〕（原-11）、アインシュタインの係数 $B$（原-4 ⑫）、バックリング $B^2$〔1/m$^2$〕（原-15） & 磁場（電）、四端子定数（電-13） & 章と単位：T、MeV、1/m$^2$ \\
$M$、$m$ & 質量 $m$〔kg〕（全章。電子は $m_{\mathrm e}$）、モル質量 $M$〔g/mol〕（原-0。原-7・8 では $M_{\text{モル}}$：原-7 は磁化とぶつけないため）、磁化 $M$（原-7）、陽子の質量 $M$（原-4 ⑤）、原子の質量 $M$〔u〕（原-11 ④）、ぶつかる相手の原子核の質量 $M$（中性子を1とした比。原-13 の中性子の減速）、体の質量 $M$（原-14 例題2。計数率の $m$ と分けた）、移動面積 $M^2$（原-15）。磁気量子数 $m_l$・$m_s$（原-7）、はかった計数率 $m$〔1/s〕（原-14）、濃縮の段の数 $m$（原-15 ⑤） & 磁化（電-7）、相互インダクタンス（電-8）、モル質量（熱-3） & 章と単位 \\
$\sigma$、$\Sigma$ & 断面積〔b〕（原-13・15）、導電率〔S/m〕（原-8・9）、標準偏差・正規分布のはば（原-3・12・14）、シュテファン・ボルツマン定数（原-2）。巨視的断面積 $\Sigma$〔1/cm〕（原-13・15。和の記号 $\sum$ とは別） & 応力（力-8）、面電荷密度・導電率（電） & 単位：b、S/m、1/cm \\
$\phi$ & 中性子束〔1/(cm$^2$·s)〕（原-13・15）、電位〔V〕（原-9） & 電位（電）、初期位相（波-1） & 単位 \\
$\rho$ & 密度〔g/cm$^3$・kg/m$^3$〕（原-0・1・7・8・11・13）、抵抗率〔Ω·m〕（原-9。原-8 では抵抗率を記号にせず、ことばで書く）、電荷密度 $\rho_{\mathrm q}$（原-9）、反応度（原-15。密度を記号にしない） & 密度（力）、電荷密度・抵抗率（電） & 単位。反応度は単位なし \\
$D$ & 吸収線量〔Gy〕（原-14）、拡散係数 $D_{\mathrm n}$・$D_{\mathrm p}$〔cm$^2$/s〕（原-9）、中性子の拡散係数 $D$〔cm〕（原-15） & 電束密度（電-3）、直径（力） & 単位：Gy、cm$^2$/s、cm \\
$N$、$n$ & 個数 $N$（原子の数：原-0・12・16、光子の数：原-2、電子の数：原-8、トランジスタの数：原-10、数えた数：原-14）、数密度 $N$（原-13〜15）、中性子数 $N$（原-1・11）、ドナー・アクセプターの密度 $N_{\mathrm d}$・$N_{\mathrm a}$、有効状態密度 $N_{\mathrm c}$・$N_{\mathrm v}$（原-9）。主量子数 $n$（原-4・5・7）、定在波・振動の量子数 $n$（原-2・6）、金の原子の数密度 $n$（原-1）、電子の数密度 $n$（原-8・9）、原子・電子の数密度 $n_{\text{原子}}$・$n_{\text{電子}}$（原-7 の例題。主量子数と分けた）、ベーテの式の電子の数密度 $n_{\mathrm e}$（原-13 ⑦）、ブラッグの条件の反射の次数 $n$（原-3・8 ②。数密度の $n$ とは別）、計数 $n$（原-12）、本当の計数率 $n$〔1/s〕（原-14）、中性子の数 $n$（原-15。燃料の数密度は $N$）。アボガドロ定数は $N_{\mathrm A}$ & 物質量 $n$（熱-3）、屈折率 $n$（波-12・電-15）、巻数 $n$（電-6） & 部と単位：1/m$^3$、cm$^{-3}$ \\
$\beta$、$\gamma$ & 速さの比 $\beta=v/c$（原-11・13）・ローレンツ因子 $\gamma$（原-11）、オペアンプの帰還率 $\beta$（原-10 ⑧）、遅発中性子の割合 $\beta$（原-15・16）、分裂でできる割合（収率）$\gamma_I$・$\gamma_X$（原-15 ⑩）、陽子の磁気回転比 $\gamma$〔MHz/T〕（原-7）。β線・γ線はことばで書き、γ線のエネルギーは $E_\gamma$（原-14） & 比熱比 $\gamma$（熱-3）、伝搬定数 $\gamma$（電-13） & 章で見分ける。磁気回転比のほかは単位なし \\
$\tau$ & 平均寿命（原-3 ⑧・原-12）、分裂からの時間（原-16）、エネルギーの閉じこめ時間 $\tau_E$（原-15 ⑫）、電子が衝突する間の平均の時間（原-8）、少数キャリアの寿命 $\tau_{\mathrm n}$・$\tau_{\mathrm p}$（原-9）、時定数（原-10）、不感時間（原-14） & 時定数（電-9・6-4）、せん断応力（力-8） & 単位はどれも s \\
$T$ & 温度〔K〕（全章。原-15・16 では温度の差を $\Delta\theta$ と書く）、半減期 $T_{1/2}$（原-12〜16）、運転した時間 $T_0$（原-15・16）、周期（原-5 ⑩）、炉の周期 $T$〔s〕（原-15 ⑥⑧）。透過率は書体を変えて $\mathcal T$（原-6） & 絶対温度（熱）、周期（力・電）、トルク（電-12） & 添字と単位：K、s \\
$R$ & リュードベリ定数 $R$・$R_\infty$〔1/m〕（原-4）、原子核の半径 $R$（原-1・11）、電子線回折のリングの半径 $R$（原-3 ④）、飛程 $R$〔cm〕（原-13）、球の炉心の半径 $R$（原-15）、気体定数（原-6・8）、抵抗〔Ω〕（原-10）。量子数の空間の球の半径は $r_{\mathrm F}$（原-8 ④。気体定数と分けた）。反射率は書体を変えて $\mathcal R$（原-6） & 抵抗（電）、気体定数（熱-3） & 単位：1/m、m、cm、Ω、J/(mol·K) \\
$W$、$w$ & 仕事関数〔eV〕（原-2・6）、空乏層の幅（原-9）、チャネルの幅（原-10）。放射線加重係数 $w_{\mathrm R}$・組織加重係数 $w_{\mathrm T}$（原-14） & 仕事（力・熱） & 単位：eV、m。$w$ は単位なし \\
$Q$ & 核反応・崩壊の Q 値〔MeV〕（原-11・12）、電気量 $Q$〔C〕（原-1 の例題）、核融合の利得 $Q$（原-15 ⑫。核融合の出力÷温めるのに入れた出力） & 熱量（熱）、電荷（電）、共振の鋭さの Q 値（波-2・電-10） & 単位：MeV、C。利得は単位なし \\
$p$ & 運動量（原-2・3・11）、正孔の密度（原-9）、4因子の共鳴を逃れる確率 $p$（原-15）。陽子の記号 p は立てた字 & 運動量（力-3）、圧力（熱） & 単位：kg·m/s、cm$^{-3}$ \\
$S$ & 阻止能〔MeV/cm〕（原-13）、面積（原-9・10）、点線源から1秒に出る光子の数 $S$（原-14） & エントロピー（熱-6）、面積（電） & 単位：MeV/cm、m$^2$ \\
$Z$ & 原子番号（全章）、有効核電荷 $Z_{\mathrm{eff}}$（原-7） & インピーダンス（電）、圧縮係数（熱-3）、分配関数（熱-15） & 単位なし \\
$a$ & 格子定数（原-8）、原子の間かく・1辺（原-0・3）、ボーア半径 $a_0$（原-4・7）、壁の厚さ（原-6）、油滴の半径（原-1）、検出器の窓の半径（原-14）、平らな炉心の厚さ（原-15）。差分法の係数 $a=\hbar^2/(2md^2)$〔eV〕（原-5 ⑫） & 加速度（力） & 単位：m。原-5 ⑫は eV \\
$L$、$\ell$ & 箱の幅・1辺（原-2・5・8）、膜や鏡までの距離（原-3・11）、拡散距離 $L$（原-9・15）、チャネルの長さ（原-10）、蒸発熱（原-16）。平均自由行程 $\ell$（原-8）、中性子の世代の時間 $\ell$（原-15） & インダクタンス（電）、潜熱（熱-1）、長さ $l$（力） & 単位：m、J/kg、s \\
$I$ & 光・放射線の強さ（原-2・13）、電流（原-9・10）、慣性モーメント（原-7）、ヨウ素135 の数 $I$〔1/cm$^3$〕（原-15 ⑩。添字の $I$・$X$ はヨウ素とキセノン）。平均励起エネルギーは $I_{\mathrm{ex}}$〔eV〕（原-13 ⑦。強さと分けた） & 電流（電）、慣性モーメント（力-6） & 単位 \\
$U$、$u$ & 位置エネルギー（ポテンシャル）$U$（原-5〜7）、チップのエネルギー $U$（原-10）、光のエネルギー密度 $u_f$（原-2・4）。原子質量単位 u は立てた字 & 内部エネルギー（熱-2）、位置エネルギー（力-4） & 単位：J・eV、u \\
$\alpha$ & 1回の正面衝突で残る割合（原-13）、レンズ・円すいの開きの半分の角（原-3 ⑥・原-14）、微細構造定数 $\alpha\approx1/137$（原-12 ⑩）、切りかわるトランジスタの割合（原-10）、温度係数 $\alpha_T$〔1/K〕（原-15 ⑩）。α線はことばで書く & 線膨張係数（力-8）、抵抗の温度係数（電-4） & 単位なし（角は rad）。$\alpha_T$ は 1/K \\
$d$、$\Delta x$ & 面・溝・原子の列の間かく $d$（原-3・8）、いちばん近づく距離 $d$（原-1）、STM のすき間 $d$（原-6）、差分法の点の間かく $d$（原-5 ⑫。不確かさ $\Delta x$ と分けた）、線源までの距離 $d$（原-14）。位置の不確かさ $\Delta x$（原-3 ⑦・原-5 ⑨。標準偏差）、二重スリットのしまの間かく $\Delta x$（原-3 ⑤。波-13 の式）、うすい層の厚さ $\Delta x$（原-13・15）。$3d$ の $d$ は軌道の名前（$s$・$p$・$d$・$f$。原-7・8） & 格子の間かく $d$・しまの間かく $\Delta x$（波-13）、距離・直径 $d$（力） & どれも長さ〔m〕。章と式で見分ける \\
$\theta$ & 角（多くの章。原-3 ③の入射の角は面に垂直から、原-3 ④・原-8 のブラッグの角は結晶の面からはかる）、デバイ温度 $\theta_{\mathrm D}$〔K〕（原-8）、温度の上がり $\Delta\theta$〔K〕（原-15・16） & 角（力・波） & 添字と単位：rad・°、K \\
$\varepsilon$、$\eta$ & 真空の誘電率 $\varepsilon_0$（原-1・9）、『熱』の1個の粒子のエネルギー $\varepsilon$・エネルギーの粒 $\Delta\varepsilon$（熱-15 を引く所。原-2・6〜8）、4因子の高速核分裂因子 $\varepsilon$・再生率 $\eta$（原-15）、空気の粘度 $\eta$（原-1）、ガモフ因子 $2\pi\eta$（原-12 ⑩）、量子効率 $\eta$（原-9） & ひずみ（力-8）、効率 $\eta$（熱・電） & 単位：F/m、Pa·s。4因子は単位なし \\
$i$、$j$ & 虚数単位 $i$（原-5〜8）と $j$（原-9・10）、電流の瞬時値 $i$（原-10） & 虚数単位 $j$（数学・電） & 章で見分ける \\
$g$、$H$ & g 因子 $g$（原-7）、相互コンダクタンス $g_{\mathrm m}$（原-10）、状態密度 $g_{\mathrm c}(E)$（原-9）、重力加速度 $g$（原-1）、中性子の世代の数 $g$（原-15）。等価線量 $H_{\mathrm T}$〔Sv〕（原-14）、差分法の行列 $H$（原-5 ⑫）、円柱の炉心の高さ $H$（原-15） & 重力加速度（力）。磁場の強さ $H$（電-7）、エンタルピー（熱-4） & 添字と単位 \\
\end{longtblr}\par}

\phbkhead{この巻の約束（記号と向き）}
{\small
\begin{itemize}[leftmargin=1.2\zw, itemsep=2pt, topsep=1pt, label={\color{bkVol}●}]
\item \textbf{光の振動数}：$f$ と書き、光子のエネルギーは $E=hf$ です。高校の教科書の $E=h\nu$ と同じ式です（原-2 ③）。$\nu$ は1回の分裂で出る中性子の数に使います（原-15）。
\item \textbf{虚数単位}：量子力学の章（原-5〜8）は、専門書と同じ $i$ と時間の部分 $e^{-iEt/\hbar}$ を使います。数学の本・『電磁気学』の $j$ と同じ数です（原-5 ④）。回路の章（原-9・10）は『電磁気学』と同じ $j$ です。
\item \textbf{部で分ける記号}：$\lambda$ は第1部と原-8・原-11 では波長、第4部（原-12〜15）では崩壊定数です。$A$ は原-1・11 では質量数、原-12〜14 では放射能です。原-14 では $E$ を実効線量に使い、エネルギーは $E_\gamma$ と添字で書きます。
\item \textbf{原子核の書き方}：${}^{A}_{Z}\mathrm X$ の左上が質量数、左下が原子番号です。本文では「セシウム137」「ウラン235」と名前と質量数で書きます。
\item \textbf{半導体}：中の電場は $\mathcal E$、ドナー・アクセプターの密度は $N_{\mathrm d}$・$N_{\mathrm a}$（アボガドロ定数 $N_{\mathrm A}$ と分けた）、熱の電圧は $k_{\mathrm B}T/e$ のまま書きます。数密度は本の慣用の cm$^{-3}$ でも書きます（$\qty{1}{cm^{-3}}=\qty{e6}{m^{-3}}$）。
\item \textbf{放射線の量}：吸収線量（Gy）は物理量、等価線量・実効線量（Sv）は防護量、測定器が示す 1\,cm 線量当量（Sv）は実用量です（原-14）。
\item \textbf{本文の丸め}：$h=\qty{6.63e-34}{J.s}$、$hc=\qty{1240}{eV.nm}$、\qty{300}{K}の $k_{\mathrm B}T=\qty{0.0259}{eV}$、$\qty{1}{u}=\qty{931.5}{MeV}$、$m_{\mathrm e}c^2=\qty{0.511}{MeV}$。
\end{itemize}\par}

%% ---------------- 核種のデータ・物性値・定数 ----------------
\bkbacktitle{核種のデータ・物性値・定数}
{\small この巻で使った値です。本文の例題の値は、この表と同じ丸めです。核データ・物性値は、はかり方や評価の年で少し変わるので、設計では最新の評価ずみのデータ（核図表・便覧）を使います。出典の欄の「原-4」などは、その値を使った章です。\phyk は、一次資料でも、それを引く2つ以上の資料でも、まだ確かめきれていない値です。\par}

\phbkhead{1　この巻で使う定数}
{\footnotesize
\begin{longtblr}{colspec={Q[l,m,wd=44mm] Q[l,m,wd=50mm] X[l,m]}, width=\linewidth, rowhead=1, row{1}={bg=bkVol!14, font=\sffamily\bfseries\footnotesize}, hlines={0.4pt,bkLine}, rowsep=1.4pt, colsep=3.5pt}
量 & 値 & 本文の丸め・章 \\
プランク定数 $h$ & \qty{6.62607015e-34}{J.s}（決めた値） & $6.63\times10^{-34}$（原-2） \\
ディラック定数 $\hbar=h/(2\pi)$ & \qty{1.054571817e-34}{J.s} & $1.05\times10^{-34}$（原-3） \\
電気素量 $e$ & \qty{1.602176634e-19}{C}（決めた値） & $1.60\times10^{-19}$（原-1） \\
光速 $c$ & \qty{299792458}{m/s}（決めた値） & $3.00\times10^{8}$ \\
ボルツマン定数 $k_{\mathrm B}$ & \qty{1.380649e-23}{J/K}（決めた値） & \qty{300}{K}で $k_{\mathrm B}T=\qty{0.02585}{eV}$。本文は 0.0259（原-8・9） \\
アボガドロ定数 $N_{\mathrm A}$ & \qty{6.02214076e23}{mol^{-1}}（決めた値） & $6.02\times10^{23}$（原-0・1） \\
電子の質量 $m_{\mathrm e}$ & \qty{9.1093837139e-31}{kg}（\qty{0.51099895}{MeV}） & \qty{0.511}{MeV}（原-1・11） \\
陽子の質量 $m_{\mathrm p}$ & \qty{1.67262192595e-27}{kg}（\qty{938.272089}{MeV}） & 原-1・11 \\
中性子の質量 $m_{\mathrm n}$ & \qty{1.67492750056e-27}{kg}（\qty{939.565422}{MeV}） & 原-11・15 \\
原子質量単位 u & \qty{1.66053906892e-27}{kg}＝\qty{931.494104}{MeV} & \qty{931.5}{MeV}（原-11） \\
電子の比電荷 $e/m_{\mathrm e}$ & \qty{1.75882000838e11}{C/kg} & $1.76\times10^{11}$（原-1） \\
真空の誘電率 $\varepsilon_0$ & \qty{8.8541878188e-12}{F/m} & $ke^2=\qty{1.440}{MeV.fm}$（原-1・11） \\
ボーア半径 $a_0$ & \qty{5.29177210544e-11}{m} & \qty{0.0529}{nm}（原-4・7） \\
リュードベリ定数 $R_\infty$ & \qty{10973731.568157}{m^{-1}} & $R_\infty hc=\qty{13.605693}{eV}$（原-4） \\
微細構造定数 $\alpha$ & $7.2973525643\times10^{-3}$（$1/137.035999$） & 原-4 \\
ボーア磁子 $\mu_{\mathrm B}$ & \qty{9.2740100657e-24}{J/T}（\qty{5.7883818e-5}{eV/T}） & 原-7 \\
電子の g 因子の大きさ & 2.00231930436 & 原-7 \\
陽子の磁気回転比 $\div2\pi$ & \qty{42.577478461}{MHz/T} & \qty{1.5}{T}で\qty{63.9}{MHz}（原-7） \\
電子のコンプトン波長 $h/(m_{\mathrm e}c)$ & \qty{2.42631023538}{pm} & 原-2・11 \\
$hc$・$\hbar c$ & \qty{1239.84198}{eV.nm}・\qty{197.3269804}{MeV.fm} & \qty{1240}{eV.nm}（全章） \\
陽子の電荷の半径 & \qty{0.84075}{fm} & 原-4 \\
シュテファン・ボルツマン定数 $\sigma$ & \qty{5.670374419e-8}{W/(m^2.K^4)} & 原-2 \\
ウィーンの定数 $b_{\mathrm W}$ & \qty{2.897771955e-3}{m.K} & 原-2 \\
放射の定数 $c_1$・$c_2$ & \qty{3.741771852e-16}{W.m^2}・\qty{1.438776877e-2}{m.K} & 原-2（熱-13） \\
ファラデー定数 $F$ & \qty{96485.33212}{C/mol} & 原-1 \\
ケイ素の格子定数 & \qty{0.5431020511}{nm} & 原-0・8 \\
\end{longtblr}\par}
{\footnotesize\color{bkGray}定数は NIST の CODATA 2022 の値（physics.nist.gov）。$\hbar c$・$hc$・$ke^2$・$k_{\mathrm B}T$ は、この表の値から計算し直した値。2019年の SI の改定で、$h$・$e$・$k_{\mathrm B}$・$N_{\mathrm A}$・$c$ は決めた値（誤差のない値）になりました。\par}

\phbkhead{2　原子の質量と結合エネルギー}
{\footnotesize
\begin{longtblr}{colspec={Q[l,m,wd=24mm] Q[c,m,wd=34mm] Q[c,m,wd=26mm] X[l,m]}, width=\linewidth, rowhead=1, row{1}={bg=bkVol!14, font=\sffamily\bfseries\footnotesize}, hlines={0.4pt,bkLine}, rowsep=1.2pt, colsep=3.5pt}
核種 & 原子の質量〔u〕 & {1核子あたりの\\結合エネルギー〔MeV〕} & 存在比・章 \\
中性子 n & 1.00866491606 & ― & CODATA 2022（原-11・15） \\
水素1 ${}^{1}$H & 1.00782503223 & ― & 99.9885\,\%（原-11） \\
水素2（重水素）${}^{2}$H & 2.01410177812 & 1.112 & 0.0115\,\%（原-11・15） \\
水素3（三重水素）${}^{3}$H & 3.0160492779 & 2.827 & 原-11・15 \\
ヘリウム3 ${}^{3}$He & 3.0160293201 & 2.573 & 原-11 \\
ヘリウム4 ${}^{4}$He & 4.00260325413 & 7.074 & 原-11・12 \\
リチウム7 ${}^{7}$Li & 7.0160034366 & 5.606 & 原-11 \\
ベリリウム9 ${}^{9}$Be & 9.012183065 & 6.463 & 原-11 \\
炭素12 ${}^{12}$C & 12（決めた値） & 7.680 & 98.93\,\%（原-11） \\
炭素14 ${}^{14}$C & 14.0032419884 & 7.520 & 原-12 \\
窒素14 ${}^{14}$N & 14.00307400443 & 7.476 & 原-11・12 \\
酸素16 ${}^{16}$O & 15.99491461957 & 7.976 & 原-11 \\
鉄56 ${}^{56}$Fe & 55.93493633 & 8.790 & 原-11 \\
ニッケル62 ${}^{62}$Ni & 61.92834537 & 8.795 & いちばん大きい（原-11） \\
クリプトン92 ${}^{92}$Kr & 91.9261731 & 8.513 & 分裂の破片（原-11・15） \\
バリウム141 ${}^{141}$Ba & 140.9144033 & 8.326 & 分裂の破片（原-11・15） \\
ポロニウム210 ${}^{210}$Po & 209.9828741 & 7.834 & 原-12 \\
ラドン222 ${}^{222}$Rn & 222.0175782 & 7.694 & 原-12 \\
ラジウム226 ${}^{226}$Ra & 226.0254103 & 7.662 & 原-12 \\
ウラン235 ${}^{235}$U & 235.0439301 & 7.591 & 0.7204\,\%（原-11・15） \\
ウラン238 ${}^{238}$U & 238.0507884 & 7.570 & 99.2742\,\%（原-11・12） \\
プルトニウム239 ${}^{239}$Pu & 239.0521636 & 7.560 & 原-15・16 \\
\end{longtblr}\par}
{\footnotesize\color{bkGray}原子の質量は、電子もふくめた中性の原子の質量（NIST "Atomic Weights and Isotopic Compositions"、AME2020 による）。結合エネルギーは $B=\bigl(Z\,m({}^{1}\mathrm H)+N\,m_{\mathrm n}-m\bigr)c^2$ で計算した（原-11 ④）。原子の質量から計算した Q 値：ラジウム226 の α崩壊 \qty{4.871}{MeV}、ポロニウム210 \qty{5.407}{MeV}、ウラン238 \qty{4.270}{MeV}、ポロニウム212 \qty{8.954}{MeV}（評価ずみの核データ ENSDF と合う。原-12）、中性子の β崩壊 \qty{0.782}{MeV}、D＋T の核融合 \qty{17.59}{MeV}（原-11・15）。\par}

\phbkhead{3　半減期と放射線のエネルギー}
{\footnotesize
\begin{longtblr}{colspec={Q[l,m,wd=30mm] Q[l,m,wd=28mm] Q[l,m,wd=40mm] X[l,m]}, width=\linewidth, rowhead=1, row{1}={bg=bkVol!14, font=\sffamily\bfseries\footnotesize}, hlines={0.4pt,bkLine}, rowsep=1.2pt, colsep=3.5pt}
核種 & 半減期 & こわれ方と主な放射線 & 出典・章 \\
中性子 & 877.75\,s（容器の方法）、887.7\,s（ビームの方法） & β崩壊して陽子に（Q 値 \qty{0.782}{MeV}） & ア（原-12） \\
水素3（トリチウム） & 12.32 年（約12.3年） & β線 & イ・ウ（原-14・16） \\
炭素14 & 5700±30 年（計算は 5730 年） & β線（窒素14 へ） & イ（原-12。数学 6-4） \\
カリウム40 & $1.248\times10^{9}$ 年（約13億年） & β線ほか。カリウムの 0.0117\,\% & イ・ウ（原-12・14） \\
コバルト60 & 5.271 年（1925.28 日） & β線、γ線 \qty{1.173}{MeV}・\qty{1.332}{MeV} & イ（原-13） \\
ストロンチウム90 & 28.91 年 & β線（イットリウム90 の β線は最大 \qty{2.28}{MeV}） & イ（原-13・16） \\
モリブデン99 & 65.924 時間（約66時間） & β線（テクネチウム99m へ） & イ（原-12。数学 6-14） \\
テクネチウム99m & 6.0072 時間 & γ線（核医学の検査） & イ（原-12。数学 3-11） \\
ヨウ素131 & 8.0252 日（約8日） & β線・γ線 & イ・ウ（原-12・14・16） \\
ヨウ素135 & 6.58 時間 & β線（キセノン135 へ） & イ（原-15） \\
キセノン135 & 9.14 時間 & β線。熱中性子を強く吸う & イ（原-15） \\
セシウム134 & 2.0652 年（約2年） & β線・γ線 & イ・ウ（原-14・16） \\
セシウム137 & 30.08 年（約30年） & β線。バリウム137m から γ線 \qty{0.662}{MeV}（85.1\,\%） & イ・ウ（原-2・12〜14・16） \\
ラドン222 & 3.8222 日（約3.8日） & α線 & イ・ウ（原-12・14） \\
ポロニウム210 & 138.376 日（約138日） & α線 \qty{5.304}{MeV} & イ・ウ（原-12・13） \\
ポロニウム212 & 294.3\,ns & α線（Q 値 \qty{8.954}{MeV}） & イ（原-12） \\
ラジウム226 & 1600 年 & α線 \qty{4.784}{MeV}（94\,\%） & イ（原-12） \\
ウラン235 & $7.04\times10^{8}$ 年 & α線 & イ（原-12・15） \\
ウラン238 & $4.468\times10^{9}$ 年（約45億年） & α線 & イ・ウ（原-12） \\
プルトニウム239 & 24110 年（約2万4千年） & α線 & イ（原-16） \\
ミューオン & \qty{2.1969803}{\micro s}（平均寿命） & 電子とニュートリノに & エ（原-11） \\
\end{longtblr}\par}
{\footnotesize\color{bkGray}出典　ア：UCNτ（arXiv:2106.10375）と Yue ら（arXiv:1309.2623）。2つのはかり方の食い違いは、まだわかっていない（原-12 ⑪）。イ：評価ずみの核データ ENSDF（IAEA LiveChart、2026年11月に読んだ）。ウ：環境省「放射線による健康影響等に関する統一的な基礎資料（令和7年度版）」の丸めた値（かっこの中）。エ：MuLan（arXiv:1211.0960）。炭素14 の計算は数学 6-4 と同じ 5730 年（平均寿命 約8270年）で、いまの評価 5700 年との差は年代で約0.5\,\%（原-12 ⑦）。ラジウム226 の 1\,g の放射能は約 $\qty{3.7e10}{Bq}$（昔の単位の 1\,Ci。原-12 ④）。\par}

\phbkhead{4　原子と分子が出す光}
{\footnotesize
\begin{longtblr}{colspec={Q[l,m,wd=40mm] Q[l,m,wd=50mm] X[l,m]}, width=\linewidth, rowhead=1, row{1}={bg=bkVol!14, font=\sffamily\bfseries\footnotesize}, hlines={0.4pt,bkLine}, rowsep=1.2pt, colsep=3.5pt}
もの & 値 & 出典・章 \\
水素の線（空気中の波長） & 656.285・486.136・434.046・410.174\,nm（バルマー系列） & オ（原-4） \\
水素のライマン α（真空中） & \qty{121.567}{nm} & オ（原-4） \\
ナトリウムの D 線 & 588.995・\qty{589.592}{nm}（3p の寿命 \qty{16.2}{ns}） & オ（原-3・4） \\
炎色反応の線 & リチウム 670.776・670.791、カリウム 766.490、カルシウム 422.673、ストロンチウム 460.733、バリウム 553.548、銅 \qty{521.820}{nm} & オ（原-4） \\
水銀（蛍光灯）・ネオン（He-Ne レーザー） & 水銀 435.833・\qty{546.074}{nm}、ネオン \qty{632.816}{nm} & オ（原-4） \\
電離エネルギー & 水素 13.598、ヘリウム 24.587（He$^+$ 54.418）、リチウム 5.392、ネオン 21.565、ナトリウム 5.139、アルゴン 15.760、カリウム \qty{4.341}{eV} & カ（原-4・7） \\
特性X線 K$\alpha_1$ & 鉄 6.404、銅 8.048、モリブデン 17.479、銀 22.163、タングステン \qty{59.319}{keV}（銅 K$\beta$ \qty{8.905}{keV}） & キ（原-4） \\
特性X線の波長 & 銅 K$\alpha_1$ \qty{0.15406}{nm}、モリブデン K$\alpha_1$ \qty{70.93}{pm} & キ（原-2・8） \\
CO（一酸化炭素） & 振動 \qty{2143.24}{cm^{-1}}、原子の間 \qty{112.83}{pm}、回転の線（$J=0\to1$）\qty{115.27}{GHz} & ク（原-6・7） \\
HCl・N$_2$ の振動 & HCl \qty{2885.31}{cm^{-1}}（原子の間 \qty{127.46}{pm}）、N$_2$ \qty{2329.92}{cm^{-1}} & ク（原-6） \\
水素分子 H$_2$ の結合 & \qty{4.478}{eV}（$\qty{36118}{cm^{-1}}$） & ク（原-7） \\
アンモニアの反転の線 & \qty{23.870}{GHz}（メーザー） & ケ（原-4・6） \\
\end{longtblr}\par}
{\footnotesize\color{bkGray}出典　オ：NIST "Handbook of Basic Atomic Spectroscopic Data"。カ：NIST Atomic Spectra Database（電離エネルギー）。キ：NIST "X-ray Transition Energies"。波長は $hc/E$ で計算した。ク：NIST Chemistry WebBook（Huber と Herzberg の2原子分子の定数）。振動の線は $\omega_e-2\omega_ex_e$ で計算した（0→1 の吸収）。ケ：米国標準局の論文と京都大学の論文で照らし合わせた値。原子の線の波長は、空気中の値（真空中より約 0.03\,\% 短い）。\par}

\phbkhead{5　固体と半導体（\qty{300}{K}）}
{\footnotesize
\begin{longtblr}{colspec={Q[l,m,wd=40mm] Q[l,m,wd=52mm] X[l,m]}, width=\linewidth, rowhead=1, row{1}={bg=bkVol!14, font=\sffamily\bfseries\footnotesize}, hlines={0.4pt,bkLine}, rowsep=1.2pt, colsep=3.5pt}
もの & 値 & 出典・章 \\
密度〔g/cm$^3$〕 & ケイ素 2.330、銅 8.960、銀 10.50、金 19.32、アルミニウム 2.699、ナトリウム 0.971、鉄 7.874、鉛 11.35 & コ（原-0・1・7・8・13） \\
ケイ素の禁制帯幅 $E_{\mathrm g}$ & \qty{1.12}{eV}（吸収のはし 約\qty{1110}{nm}） & サ（原-2・8・9）\phyk \\
ケイ素の有効状態密度 $N_{\mathrm c}$・$N_{\mathrm v}$ & $2.8\times10^{19}$・$\qty{1.04e19}{cm^{-3}}$ & サ（原-9）\phyk \\
ケイ素の真性キャリア密度 $n_{\mathrm i}$ & 本により $6.7\times10^{9}$〜$\qty{1.45e10}{cm^{-3}}$。例題は $1.0\times10^{10}$ & サ（原-9）\phyk \\
ケイ素の移動度 $\mu_{\mathrm n}$・$\mu_{\mathrm p}$ & 1350・\qty{480}{cm^2/(V.s)} & サ（原-9・10）\phyk \\
ケイ素の比誘電率 & 11.7 & サ（原-9）\phyk \\
ほかの禁制帯幅 & ゲルマニウム 0.66、GaAs 1.42、GaN 約3.4、石英ガラス 約 \qty{9}{eV}（8〜10\,eV） & サ・シ（原-2・8・9）\phyk \\
GaAs の電子の有効質量 & 約 $0.067\,m_{\mathrm e}$ & サ（原-5） \\
銅の自由電子 & 数密度 $\qty{8.49e28}{m^{-3}}$、$E_{\mathrm F}=\qty{7.0}{eV}$、$v_{\mathrm F}=\qty{1.57e6}{m/s}$、衝突の間の時間 \qty{2.4e-14}{s} & ス（原-7・8。電-4） \\
銅のデバイ温度 & \qty{343}{K} & セ（原-8）\phyk \\
仕事関数（目安） & セシウム 約2.1、カリウム 約2.3、タングステン 約4.5、銅 約\qty{4.7}{eV}（表面の状態で 0.1〜\qty{0.5}{eV} 違う） & ソ（原-2・6） \\
超伝導の臨界温度 & ニオブ・チタン 約\qty{9.5}{K}、液体窒素の沸点 \qty{77.3}{K} & セ・ク（原-8）\phyk \\
鉄の1原子あたりの磁石の強さ & 約 $2.2\,\mu_{\mathrm B}$ & セ（原-7）\phyk \\
\end{longtblr}\par}
{\footnotesize\color{bkGray}出典　コ：NIST "X-Ray Mass Attenuation Coefficients" の表1。ケイ素の密度は格子定数から計算すると 2.329\,g/cm$^3$（原-8 ③）。サ：S.\,M.\,Sze の半導体の教科書の表を引く資料2つ以上で照らし合わせた値（一次資料のデータベースに届かなかった）。シ：石英ガラスの反射率・光伝導の論文。ス：電-4 の値（密度と1\,mol の質量から1原子に1個として計算）と、それから計算した値。セ：固体物理の教科書の値、製品の資料。ソ：大学の解説の表と Michaelson（1977）の値で照らし合わせた目安。\par}

\phbkhead{6　放射線と物質}
{\footnotesize
\begin{longtblr}{colspec={Q[l,m,wd=26mm] Q[c,m,wd=17mm] Q[c,m,wd=20mm] Q[c,m,wd=20mm] Q[c,m,wd=20mm] X[l,m]}, width=\linewidth, rowhead=1, row{1}={bg=bkVol!14, font=\sffamily\bfseries\footnotesize}, hlines={0.4pt,bkLine}, rowsep=1.2pt, colsep=3pt}
物質 & {密度\\〔g/cm$^3$〕} & {$\mu/\rho$\\60\,keV} & {$\mu/\rho$\\0.662\,MeV} & {$\mu/\rho$\\1.25\,MeV} & {0.662\,MeV の\\半価層} \\
骨（緻密な骨） & 1.92 & 0.315 & 0.0797 & ― & \qty{4.5}{cm} \\
軟組織 & 1.06 & 0.205 & 0.0849 & ― & \qty{7.7}{cm} \\
水 & 1.00 & 0.206 & 0.0857 & 0.0632 & \qty{8.1}{cm} \\
コンクリート & 2.30 & 0.266 & 0.0788 & 0.0581 & \qty{3.8}{cm} \\
アルミニウム & 2.699 & 0.278 & 0.0746 & 0.0550 & \qty{3.4}{cm} \\
鉄 & 7.874 & 1.21 & 0.0735 & 0.0535 & \qty{1.2}{cm} \\
鉛 & 11.35 & 5.02 & 0.111 & 0.0588 & \qty{0.55}{cm} \\
\end{longtblr}\par}
{\footnotesize\color{bkGray}質量減弱係数 $\mu/\rho$ の単位は cm$^2$/g（NIST "X-Ray Mass Attenuation Coefficients"。0.662\,MeV は 0.6\,MeV と 0.8\,MeV の値から補間した。骨・軟組織は ICRU の標準の組成。原-13 ④）。半価層は $\ln2/\mu$ で計算した。\par}
\vspace{1mm}
{\footnotesize
\begin{longtblr}{colspec={Q[l,m,wd=44mm] Q[l,m,wd=56mm] X[l,m]}, width=\linewidth, rowhead=1, row{1}={bg=bkVol!14, font=\sffamily\bfseries\footnotesize}, hlines={0.4pt,bkLine}, rowsep=1.2pt, colsep=3.5pt}
もの & 値 & 出典・章 \\
陽子の水の中の飛程（CSDA） & 70\,MeV で \qty{4.08}{cm}、150\,MeV で \qty{15.8}{cm}、225\,MeV で \qty{31.7}{cm}、250\,MeV で \qty{37.9}{cm} & タ（原-13） \\
α線の空気の中の飛程 & \qty{5.3}{MeV}で約\qty{4.0}{cm} & タ（原-13） \\
電子のアルミニウムの中の飛程（CSDA） & 2.0\,MeV で \qty{1.22}{g/cm^2}、2.5\,MeV で \qty{1.55}{g/cm^2} & タ（原-13） \\
空気で1組のイオンを作るエネルギー $W$ & \qty{33.97}{eV}（約\qty{34}{eV}） & チ（原-13・14） \\
熱中性子 & \qty{0.0253}{eV}、速さ \qty{2200}{m/s}（\qty{20.4}{\celsius}の $k_{\mathrm B}T$） & 計算（原-3・13・15） \\
\end{longtblr}\par}
{\footnotesize\color{bkGray}出典　タ：NIST の PSTAR・ASTAR・ESTAR（水・乾いた空気・アルミニウム）。チ：ICRU Report 90 が確かめた値（それを引く論文2つで照らし合わせた）。\par}

\phbkhead{7　中性子と原子炉}
{\footnotesize
\begin{longtblr}{colspec={Q[l,m,wd=50mm] Q[l,m,wd=48mm] X[l,m]}, width=\linewidth, rowhead=1, row{1}={bg=bkVol!14, font=\sffamily\bfseries\footnotesize}, hlines={0.4pt,bkLine}, rowsep=1.2pt, colsep=3.5pt}
もの & 値 & 出典・章 \\
ウラン235 の分裂の断面積（熱中性子） & \qty{585}{b}（標準 587.3±1.4\,b） & ツ・テ（原-13・15） \\
ウラン235 の捕獲の断面積（熱中性子） & \qty{99}{b}（標準 99.5±1.3\,b） & ツ・テ（原-15） \\
ウラン238 の捕獲の断面積 & \qty{2.68}{b} & テ・ト（原-15） \\
水素1 の吸収の断面積 & \qty{0.3326}{b} & ト（原-13・15） \\
コバルト59 の吸収の断面積 & \qty{37.18}{b} & ト（原-13） \\
キセノン135 の捕獲の断面積 & 約 $\qty{2.6e6}{b}$ & テ（原-15） \\
1回の分裂で出る中性子の数 $\nu$ & 2.43（標準 2.425±0.011） & ツ（原-15） \\
遅発中性子の割合 $\beta$ & 0.0065 & ナ（原-15）\phyk \\
1回の分裂のエネルギー & \qty{202.8}{MeV}（破片の運動 166.2、即発中性子 4.8、即発γ線 8.0、破片のβ線 7.0、破片のγ線 7.2、ニュートリノ 9.6）。本文は約\qty{200}{MeV} & ニ（原-0・11・15・16） \\
分裂でできる割合（収率） & ヨウ素135 6.39\,\%、キセノン135 が直接できる分 約0.24\,\% & イ（原-15） \\
崩壊熱の目安の式 & {$P/P_0\approx$\\ \quad$0.066\,[t^{-0.2}-(t+T_0)^{-0.2}]$\\（$t$・$T_0$ は秒）。3年運転して止めた1日後 約0.5\,\%、1年後 約0.05\,\%} &ヌ（原-15・16）\phyk \\
\end{longtblr}\par}
{\footnotesize\color{bkGray}出典　ツ：IAEA Neutron Data Standards 2017 の表11（熱中性子定数）。本文の 585\,b・99\,b は評価ずみのライブラリ（JENDL-5・ENDF/B-VIII など）の値で、標準と 0.5\,\% 以内。テ：Pritychenko（arXiv:2503.18990）の評価ずみライブラリの表。ト：NIST NCNR の中性子の散乱長と断面積の表（Sears 1992）。ナ：米国エネルギー省の手引き（DOE-HDBK-1019）と arXiv:2101.09338 で照らし合わせた（$\nu_{\mathrm d}/\nu=0.0157/2.425=0.0065$）。ニ：Algora ら（arXiv:2007.07918）の表。ヌ：ウェイとウィグナーの目安の式（係数は本により 0.062〜0.066）。イ：表3 と同じ。濃縮していないウランのウラン235 は 0.7204\,\%、軽水炉の燃料は 3〜5\,\%（環境省の基礎資料）。\par}

\phbkhead{8　線量と防護}
{\footnotesize
\begin{longtblr}{colspec={Q[l,m,wd=44mm] Q[l,m,wd=62mm] X[l,m]}, width=\linewidth, rowhead=1, row{1}={bg=bkVol!14, font=\sffamily\bfseries\footnotesize}, hlines={0.4pt,bkLine}, rowsep=1.2pt, colsep=3.5pt}
もの & 値 & 出典・章 \\
放射線加重係数 $w_{\mathrm R}$ & γ線・X線・β線 1、陽子 2、α線・重いイオン 20、中性子 2.5〜21（エネルギーで変わる） & ウ（原-14） \\
組織加重係数 $w_{\mathrm T}$ & 骨髄・結腸・肺・胃・乳房 各0.12、生殖腺 0.08、ぼうこう・食道・肝臓・甲状腺 各0.04、骨の表面・脳・だ液腺・皮ふ 各0.01、残りの組織の合計 0.12（全部で1） & ウ（原-14） \\
経口の預託実効線量係数（大人） & ストロンチウム90 0.028、ヨウ素131 0.022、セシウム134 0.019、セシウム137 0.013、プルトニウム239 0.25、トリチウム \qty{0.000018}{\micro Sv/Bq} & ウ（原-14・16） \\
自然放射線（1年） & 日本 \qty{2.1}{mSv}、世界 \qty{2.4}{mSv}。日本の内訳：宇宙線 0.3、大地 0.33、ラドン222 0.37、ラドン220 0.09、食品（鉛210・ポロニウム210）0.80、カリウム40 0.18 など & ウ（原-14） \\
医療（日本、1年の平均） & \qty{2.6}{mSv} & ウ（原-14） \\
生活の空間線量率 & 0.01〜\qty{1}{\micro Sv/h}（場所で違う） & ウ（原-14） \\
がんのリスク（防護の目安） & \qty{100}{mSv}あたり 約0.5\,\%（1000人中 約300人ががんで亡くなる所に約5人を足す） & ウ（原-14・16） \\
線量限度（一般の人） & 1年に\qty{1}{mSv}（安全と危険の境界線ではない） & ウ（原-14） \\
線量限度（放射線を扱う仕事） & 5年で\qty{100}{mSv}、1年で\qty{50}{mSv}、女性は3か月で\qty{5}{mSv}、緊急作業 \qty{100}{mSv}。管理区域は3か月で\qty{1.3}{mSv}をこえるおそれのある所 & ノ（原-14） \\
セシウム137 の 1\,cm 線量当量率定数 & \qty{0.0927}{\micro Sv.m^2/(MBq.h)}（空気カーマ率定数を計算すると 0.076\,μGy·m$^2$/(MBq·h)） & ハ（原-14）\phyk \\
\end{longtblr}\par}
{\footnotesize\color{bkGray}出典　ウ：環境省「放射線による健康影響等に関する統一的な基礎資料（令和7年度版）」（係数は ICRP 2007年勧告・Publ.\,119）。ノ：電離放射線障害防止規則 第3条・第4条・第7条（e-Gov 法令）。ハ：日本アイソトープ協会『アイソトープ手帳』の値を引く資料2つで照らし合わせた値。空気カーマ率定数は、NIST の空気の質量エネルギー吸収係数から計算した（原-14）。線量の数は、ことわりのない限り大人の値です。\par}

\Needspace{14\baselineskip}\phbkhead{9　単位と換算}
{\footnotesize
\begin{longtblr}{colspec={Q[l,m,wd=34mm] Q[l,m,wd=56mm] X[l,m]}, width=\linewidth, rowhead=1, row{1}={bg=bkVol!14, font=\sffamily\bfseries\footnotesize}, hlines={0.4pt,bkLine}, rowsep=1.2pt, colsep=3.5pt}
単位 & 換算 & 章 \\
電子ボルト eV & $\qty{1}{eV}=\qty{1.602176634e-19}{J}$。$\qty{1}{keV}=10^3$\,eV、$\qty{1}{MeV}=10^6$\,eV & 電-2・原-1 \\
光子のエネルギーと波長 & $E\,[\mathrm{eV}]=1240/\lambda\,[\mathrm{nm}]$ & 原-2 \\
温度とエネルギー & \qty{1}{eV}は約\qty{11600}{K}の $k_{\mathrm B}T$。\qty{300}{K}で\qty{0.0259}{eV} & 原-2・8 \\
原子質量単位 u（ダルトン Da） & $\qty{1}{u}=\qty{1.66054e-27}{kg}$＝\qty{931.494}{MeV}（$c^2$ をかけたエネルギー） & 原-11 \\
長さ & $\qty{1}{nm}=10^{-9}$\,m、$\qty{1}{pm}=10^{-12}$\,m、$\qty{1}{fm}=10^{-15}$\,m。専門書の $\qty{1}{\angstrom}=\qty{0.1}{nm}$ & 原-1・3 \\
バーン b（断面積） & $\qty{1}{b}=10^{-28}$\,m$^2$＝$10^{-24}$\,cm$^2$ & 原-13 \\
分光の波数 cm$^{-1}$ & 1\,cm に入る波の数。$\qty{1}{cm^{-1}}$ は \qty{29.98}{GHz}、$\qty{1.240e-4}{eV}$ & 原-6 \\
数密度 & $\qty{1}{cm^{-3}}=10^{6}$\,m$^{-3}$ & 原-9 \\
ベクレル Bq & 1秒に1個こわれる。昔の単位 $\qty{1}{Ci}=\qty{3.7e10}{Bq}$ & 原-12 \\
グレイ Gy・シーベルト Sv & $\qty{1}{Gy}=\qty{1}{J/kg}$（物理量）。Sv も J/kg だが、係数をかけた防護量 & 原-14 \\
年 & 1年＝365.25日＝$\qty{3.156e7}{s}$ & 原-12 \\
\end{longtblr}\par}
{\footnotesize\color{bkGray}接頭語：T（テラ、$10^{12}$）、G（ギガ、$10^9$）、M（メガ、$10^6$）、k（キロ、$10^3$）、m（ミリ、$10^{-3}$）、\textmu（マイクロ、$10^{-6}$）、n（ナノ、$10^{-9}$）、p（ピコ、$10^{-12}$）、f（フェムト、$10^{-15}$）。\par}

%% ---------------- なぜ？の道すじ ----------------
%% 第2巻の道すじの終点が原の章だったもの（太陽の表面の温度→原-2、ウランの濃縮→原-15）を、この巻の中の道として続けた。
%% 章の問い（もくじ）17本のあとに、章の中の身近な問い（まだわかっていないこと・やってみよう・節の題）を足した。
\bkbacktitle{なぜ？の道すじ}
{\small 気になる「なぜ？」から、読む章の順番をたどる路線図です。札の横の数はこの巻のページです。ほかの巻の札（力・熱・波・電）は、物理シリーズの『力学』『熱学』『波動』『電磁気学』の章です。使う数学の章は、下の小さな札です。\par}
\vspace{1mm}
\phwhyroute{原子1個は、どれくらい小さい？}{原-0}{原-0,原-1}{2-7}
\phwhyroute{原子の中は、ほとんどからっぽだって本当？}{原-1}{電-1,原-1}{2-7}
\phwhyroute{熱した鉄は、なぜ温度で光の色が変わる？}{原-2}{熱-13,原-2}{5-14}
\phwhyroute{電子顕微鏡は、なぜふつうの顕微鏡より小さなものが見える？}{原-3}{波-13,原-2,原-3}{4-9,9-4}
\phwhyroute{花火の色は、なぜ材料によって変わる？}{原-4}{原-2,原-4}{4-2}
\phwhyroute{電子は、原子の中のどこにいる？}{原-5}{波-6,原-3,原-5}{7-8,9-1}
\phwhyroute{USBメモリは、どうやって電子を閉じこめて記録する？}{原-6}{原-5,原-6}{5-14}
\phwhyroute{ネオンは、なぜほかの物質とほとんど結びつかない？}{原-7}{原-4,原-5,原-7}{8-12}
\phwhyroute{金属は電気を通すのに、ガラスはなぜ通さない？}{原-8}{電-4,原-5,原-8}{5-14}
\phwhyroute{LEDは、どうやって電気を直接、光に変える？}{原-9}{原-2,原-8,原-9}{3-12,5-14}
\phwhyroute{スマホのチップは、どうやって計算している？}{原-10}{原-9,原-10}{2-17}
\phwhyroute{太陽は、何をエネルギーのもとにして光っている？}{原-11}{原-11}{4-4,5-17}
\phwhyroute{大昔の木や骨の年代は、どうやってわかる？}{原-12}{原-11,原-12}{3-11,6-4}
\phwhyroute{X線写真で、なぜ骨が白くうつる？}{原-13}{原-2,原-13}{3-11,5-14}
\phwhyroute{目に見えない放射線を、どうやってはかる？}{原-14}{原-12,原-13,原-14}{3-14,3-17}
\phwhyroute{原子炉は、どうやって反応を一定に保つ？}{原-15}{原-11,原-12,原-15}{6-4,9-9}
\phwhyroute{使い終わった核燃料は、どうする？}{原-16}{原-12,原-15,原-16}{3-11}
\phwhyroute{太陽の表面の温度は、どうやってわかる？}{原-2}{熱-13,原-2}{5-14}
\phwhyroute{ウランの濃縮では、どうやって軽い原子を集める？}{原-15}{熱-14,原-15}{3-12,3-17}
\phwhyroute{電子に、大きさはある？}{原-0}{原-0,原-7}{}
\phwhyroute{陽子の大きさは、どれだけか？}{原-4}{原-1,原-4}{}
\phwhyroute{光が当たってから、電子が飛び出すまでの時間は？}{原-2}{原-2}{}
\phwhyroute{「観測すると1点に決まる」とは、どういうことか？}{原-3}{原-3,原-5}{}
\phwhyroute{量子ドットは、なぜ大きさで色が変わる？}{原-5}{原-3,原-5}{}
\phwhyroute{走査トンネル顕微鏡は、どうやって原子を見る？}{原-6}{原-5,原-6}{5-14}
\phwhyroute{MRI は、どうやって体の中を見る？}{原-7}{電-6,原-7}{}
\phwhyroute{高温超伝導のしくみは？}{原-8}{原-8}{}
\phwhyroute{太陽電池は、どうやって光から電気をつくる？}{原-9}{原-2,原-9}{3-12}
\phwhyroute{青色 LED は、なぜ大きな電流で効率が下がる？}{原-9}{原-9}{}
\phwhyroute{トランジスタは、どこまで小さくできる？}{原-10}{原-9,原-10}{}
\phwhyroute{GPS の時計は、なぜずれる？}{原-11}{力-1,原-11}{4-4,5-17}
\phwhyroute{中性子の寿命は、なぜ2通りにはかると違う？}{原-12}{原-12}{}
\phwhyroute{陽子線のがん治療は、なぜ体の深い所をねらえる？}{原-13}{原-11,原-13}{}
\phwhyroute{自然の放射線を、1年にどれくらい受けている？}{原-14}{原-12,原-14}{}
\phwhyroute{100 mSv より少ない線量で、がんは増えるか？}{原-14}{原-14}{}
\phwhyroute{原子炉を止めたあとも、なぜ冷やし続ける？}{原-15}{原-12,原-15}{6-4}
\phwhyroute{福島第一原子力発電所の事故では、何が起きた？}{原-16}{原-15,原-16}{}

%% ---------------- チェック表 ----------------
\bkbacktitle{できるようになったことチェック表}
{\small 各章の「ふりかえり」を集めた表です。章を学び終えたら1回め、シリーズを読み終えて専門の本へ進む前に2回めのチェックをしよう。自信がない所は、ページの数字の所へもどろう。\par}\vspace{2mm}
% 自動生成（tools/mkchecklist.py p5）。各章の「ふりかえり」の項目を集めたもの（第13回）
\begingroup\footnotesize
\def\bkcb{\tikz[baseline=-0.5ex]\draw[\bkTF, line width=0.7pt, rounded corners=0.5pt] (0,-0.3em) rectangle (0.8em,0.5em);}
\begin{longtblr}{colspec={Q[c,m,wd=11mm] X[l,m] Q[c,m,wd=10mm] Q[c,m,wd=9mm] Q[c,m,wd=9mm]}, width=\linewidth, rowhead=1, colsep=3pt,
  row{1}={bg=bkVol!14, font=\sffamily\bfseries\scriptsize}, hlines={0.4pt,bkLine}, rowsep=0.7pt}
章 & できるようになったこと & ページ & 1回め & 2回め \\
{\sffamily\bfseries\color{\bkTT}原-0} & 原子1個の大きさと、1円玉の中の原子の数を、けたで言える。スマホのチップと原子炉が、原子のどの部分の物理で動くかを説明できる。 & \pageref{ph:原-0:sec3} & \bkcb & \bkcb \\
 & 何を無視したかを書いてから、原子の数・1個の大きさ・1個あたりのエネルギーを、7つの段で求められる。 & \pageref{ph:原-0:ex1} & \bkcb & \bkcb \\
 & 身近な「なぜ？」の答えが、この巻のどの章にあるかを探せる。 & \pageref{ph:原-0:sec4} & \bkcb & \bkcb \\
{\sffamily\bfseries\color{\bkTT}原-1} & 陰極線が電子の流れであることを、実験の結果から説明できる。 & \pageref{ph:原-1:sec1} & \bkcb & \bkcb \\
 & トムソンの方法で、電子の比電荷を求められる。 & \pageref{ph:原-1:sec2} & \bkcb & \bkcb \\
 & ミリカンの方法で、油滴の電荷と電気素量を求められる。 & \pageref{ph:原-1:sec3} & \bkcb & \bkcb \\
 & 原子核が小さく重いことを、いちばん近づく距離で説明できる。 & \pageref{ph:原-1:sec5} & \bkcb & \bkcb \\
 & 散乱角の式で、大きく曲がる α線の割合を計算できる。 & \pageref{ph:原-1:sec8} & \bkcb & \bkcb \\
 & 質量分析で、同位体のイオンが届く場所の違いを求められる。 & \pageref{ph:原-1:sec11} & \bkcb & \bkcb \\
{\sffamily\bfseries\color{\bkTT}原-2} & 熱した物の色が温度で変わるわけを、光の分布の山の動きで説明できる。 & \pageref{ph:原-2:sec1} & \bkcb & \bkcb \\
 & 光子のエネルギー $E=hf=hc/\lambda$ を、eV と J で求められる。 & \pageref{ph:原-2:sec3} & \bkcb & \bkcb \\
 & 光電効果のきまりを光子で説明し、$K_0=hf-W$ と阻止電圧を使える。 & \pageref{ph:原-2:sec5} & \bkcb & \bkcb \\
 & 光子の運動量 $p=h/\lambda$ から、放射圧を求められる。 & \pageref{ph:原-2:sec6} & \bkcb & \bkcb \\
 & コンプトン効果の式で、散乱した X線・γ線の波長とエネルギーを求められる。 & \pageref{ph:原-2:sec7} & \bkcb & \bkcb \\
 & モードの数と平均のエネルギーから、プランクの式・ウィーンの変位則・$\sigma$ を導ける。 & \pageref{ph:原-2:sec11} & \bkcb & \bkcb \\
{\sffamily\bfseries\color{\bkTT}原-3} & 加速した電圧から、電子の波長（ド・ブロイ波長）を求められる。 & \pageref{ph:原-3:sec2} & \bkcb & \bkcb \\
 & 電子の回折と、電子1個ずつの二重スリットのしま、電子顕微鏡が細かい物を見分けられるわけを、波の式で説明できる。 & \pageref{ph:原-3:sec5} & \bkcb & \bkcb \\
 & 波の集まり（波束）のはばから不確定性関係を導き、原子の中の電子のエネルギーの大きさを見積もれる。 & \pageref{ph:原-3:sec7} & \bkcb & \bkcb \\
 & 電子の波の分散関係から、群速度が粒の速さになることを示せる。 & \pageref{ph:原-3:sec9} & \bkcb & \bkcb \\
{\sffamily\bfseries\color{\bkTT}原-4} & 炎色反応と線スペクトルを、原子のエネルギーの段（準位）の差から説明できる。 & \pageref{ph:原-4:sec3} & \bkcb & \bkcb \\
 & ボーアの模型から水素原子のエネルギー準位を導き、スペクトルの波長を求められる。 & \pageref{ph:原-4:sec5} & \bkcb & \bkcb \\
 & X線管の電圧から最短波長を、モーズリーの法則から特性X線のエネルギーを求められる。 & \pageref{ph:原-4:sec9} & \bkcb & \bkcb \\
 & 誘導放出と反転分布、アインシュタインの係数から、レーザーのしくみを説明できる。 & \pageref{ph:原-4:sec11} & \bkcb & \bkcb \\
{\sffamily\bfseries\color{\bkTT}原-5} & 箱に閉じこめた電子のエネルギー $E_n=n^2h^2/(8mL^2)$ を、定在波と $\lambda=h/p$ から求められる。 & \pageref{ph:原-5:sec3} & \bkcb & \bkcb \\
 & $|\psi|^2$ を「見つかる確率」として読み、ある範囲に電子が見つかる確率を求められる。 & \pageref{ph:原-5:sec4} & \bkcb & \bkcb \\
 & シュレディンガー方程式を立て、無限に深い井戸を境界の条件で解ける。 & \pageref{ph:原-5:sec8} & \bkcb & \bkcb \\
 & 重ね合わせと期待値を計算し、差分法で固有値を数値で求められる。 & \pageref{ph:原-5:sec12} & \bkcb & \bkcb \\
{\sffamily\bfseries\color{\bkTT}原-6} & 壁の中で波動関数が $e^{-\kappa x}$ で弱まることを使い、すりぬける割合 $e^{-2\kappa a}$ を見積もれる。 & \pageref{ph:原-6:sec3} & \bkcb & \bkcb \\
 & フラッシュメモリや走査トンネル顕微鏡のしくみを、トンネル効果で説明できる。 & \pageref{ph:原-6:sec4} & \bkcb & \bkcb \\
 & 境目のつながりから、階段の反射率・有限の井戸の準位・正確な透過率を求められる。 & \pageref{ph:原-6:sec6} & \bkcb & \bkcb \\
 & 調和振動子の準位 $(n+\tfrac12)\hbar\omega$ を使い、分子の振動の光と比熱を説明できる。 & \pageref{ph:原-6:sec8} & \bkcb & \bkcb \\
{\sffamily\bfseries\color{\bkTT}原-7} & 周期表の横の並びの数（2・8・8・18…）が、原子の中の電子の「席」の数で決まることを説明できる。 & \pageref{ph:原-7:sec1} & \bkcb & \bkcb \\
 & 量子数で席を数えて $2n^2$ を導き、ナトリウム・鉄・銅などの電子配置を書ける。 & \pageref{ph:原-7:sec4} & \bkcb & \bkcb \\
 & 水素原子の電子の雲から、電子が見つかりやすい距離や見つかる確率を求め、電離エネルギーから有効核電荷を見積もれる。 & \pageref{ph:原-7:sec3} & \bkcb & \bkcb \\
 & 磁場の中の電子や陽子の向きが2つに分かれることから、ESR や MRI の周波数を求められる。 & \pageref{ph:原-7:sec6} & \bkcb & \bkcb \\
 & 3次元のシュレディンガー方程式に水素原子の $1s$ の波動関数を代入して $-13.6$ eV を確かめ、化学結合や分子の回転のしくみを読める。 & \pageref{ph:原-7:sec10} & \bkcb & \bkcb \\
{\sffamily\bfseries\color{\bkTT}原-8} & 金属・半導体・絶縁体の違いを、電子のエネルギーの帯（バンド）と、帯のすき間で説明できる。 & \pageref{ph:原-8:sec8} & \bkcb & \bkcb \\
 & ブラッグの条件で X線の回折の角を求め、格子定数から結晶の密度や原子の数を計算できる。 & \pageref{ph:原-8:sec2} & \bkcb & \bkcb \\
 & 3次元の箱で電子の席を数えて、金属のフェルミエネルギーとフェルミ速度を求められる。 & \pageref{ph:原-8:sec4} & \bkcb & \bkcb \\
 & フェルミ分布で、電子が席をうめる割合を求め、金属の電子の比熱が小さいわけを説明できる。 & \pageref{ph:原-8:sec5} & \bkcb & \bkcb \\
 & くし形の結晶の帯をプログラムで求め、平均自由行程・ウィーデマン・フランツの法則・デバイの $T^3$ 則を使える。 & \pageref{ph:原-8:sec10} & \bkcb & \bkcb \\
{\sffamily\bfseries\color{\bkTT}原-9} & 光電池に光を当てると電気ができ、LED は決まった向きにつないだときだけ光ることを、実験で確かめられる。 & \pageref{ph:原-9:sec1} & \bkcb & \bkcb \\
 & n形・p形の半導体の電子と正孔の数と抵抗率を求め、pn 接合が電流を一方の向きにだけ通すしくみを、図で説明できる。 & \pageref{ph:原-9:sec4} & \bkcb & \bkcb \\
 & ダイオードの電流と電圧の関係（約 60 mV で電流10倍）、LED の色と禁制帯の幅の関係を使える。 & \pageref{ph:原-9:sec7} & \bkcb & \bkcb \\
 & 太陽電池の電流と電圧のグラフから、開放電圧・いちばん大きな電力・曲線因子・効率を求められる。 & \pageref{ph:原-9:sec10} & \bkcb & \bkcb \\
 & フェルミ分布の裾からキャリアの数の式を、ドリフトと拡散のつり合いから内蔵電位とダイオードの式を導ける。 & \pageref{ph:原-9:sec11} & \bkcb & \bkcb \\
{\sffamily\bfseries\color{\bkTT}原-10} & スイッチ2つで AND と OR の回路を作り、オンとオフで2進数のたし算を表せる。 & \pageref{ph:原-10:sec1} & \bkcb & \bkcb \\
 & バイポーラトランジスタと MOSFET がスイッチ・増幅器としてはたらくしくみを pn 接合と電場で説明し、相互コンダクタンス・増幅度・オン抵抗を求められる。 & \pageref{ph:原-10:sec3} & \bkcb & \bkcb \\
 & CMOS の論理ゲートのしくみを説明し、チップの消費電力と信号のおくれを計算できる。 & \pageref{ph:原-10:sec7} & \bkcb & \bkcb \\
 & 理想オペアンプの反転・非反転増幅の倍率を求め、センサーの小さな電圧を A/D 変換する流れを説明できる。 & \pageref{ph:原-10:sec8} & \bkcb & \bkcb \\
 & MOSFET の電流の式を反転層の電荷とドリフトから導き、パワー半導体の損失と、ニュートン法で回路を解く方法を使える。 & \pageref{ph:原-10:sec10} & \bkcb & \bkcb \\
{\sffamily\bfseries\color{\bkTT}原-11} & 原子番号・質量数・同位体ということばで、原子核のなかみを説明できる。 & \pageref{ph:原-11:sec3} & \bkcb & \bkcb \\
 & 原子の質量の表から、結合エネルギーと核反応の Q 値を計算できる。 & \pageref{ph:原-11:sec4} & \bkcb & \bkcb \\
 & 光の時計から時間のおくれを導き、GPS の時計のずれを計算できる。 & \pageref{ph:原-11:sec7} & \bkcb & \bkcb \\
 & $E^2=(pc)^2+(mc^2)^2$ から、$E=mc^2$ とコンプトンの式を導ける。 & \pageref{ph:原-11:sec8} & \bkcb & \bkcb \\
 & 半経験的質量公式で、結合エネルギーの曲線の形を説明できる。 & \pageref{ph:原-11:sec10} & \bkcb & \bkcb \\
{\sffamily\bfseries\color{\bkTT}原-12} & α崩壊・β崩壊・γ線で、原子核がどう変わるかを書ける。 & \pageref{ph:原-12:sec2} & \bkcb & \bkcb \\
 & 半減期を使って、残る量や、たった時間を求められる。 & \pageref{ph:原-12:sec3} & \bkcb & \bkcb \\
 & 放射能（Bq）と原子の数の関係 $A=\lambda N$ を使える。 & \pageref{ph:原-12:sec4} & \bkcb & \bkcb \\
 & 崩壊の法則を微分方程式で書き、数えるばらつきを見積もれる。 & \pageref{ph:原-12:sec5} & \bkcb & \bkcb \\
 & 親と子の崩壊（放射平衡）と、α崩壊のトンネルを式で読める。 & \pageref{ph:原-12:sec8} & \bkcb & \bkcb \\
{\sffamily\bfseries\color{\bkTT}原-13} & α線・β線・γ線（X線）・中性子線が、どんな物で止まるか（透過力）と、物を電離するしくみを説明できる。 & \pageref{ph:原-13:sec1} & \bkcb & \bkcb \\
 & 光子の減弱の式 $I=I_0e^{-\mu x}$ と半価層を使って、X線写真や遮へいの厚さを計算できる。 & \pageref{ph:原-13:sec3} & \bkcb & \bkcb \\
 & 光電効果・コンプトン散乱・電子対生成の違いと、荷電粒子・中性子の止まり方の違いを説明できる。 & \pageref{ph:原-13:sec4} & \bkcb & \bkcb \\
 & 阻止能とブラッグピーク、断面積と平均自由行程、中性子の減速を、式で読める。 & \pageref{ph:原-13:sec7} & \bkcb & \bkcb \\
 & 放射化の式 $A=N\sigma\phi(1-e^{-\lambda t})$ で、飽和放射能に近づくようすを計算できる。 & \pageref{ph:原-13:sec9} & \bkcb & \bkcb \\
{\sffamily\bfseries\color{\bkTT}原-14} & 放射線は目に見えないが、電離を使って数えられることと、Bq・Gy・Sv の意味の違いを説明できる。 & \pageref{ph:原-14:sec2} & \bkcb & \bkcb \\
 & 吸収線量から等価線量・実効線量を計算し、自然放射線や医療の線量と比べられる。 & \pageref{ph:原-14:sec5} & \bkcb & \bkcb \\
 & 点線源の逆2乗の法則と遮へいを組み合わせ、時間・距離・遮へいの3原則で線量を見積もれる。 & \pageref{ph:原-14:sec7} & \bkcb & \bkcb \\
 & 計数のばらつき（$\sqrt N$）・数え落とし・検出限界・立体角を使って、測定の結果を正しく読める。 & \pageref{ph:原-14:sec3} & \bkcb & \bkcb \\
 & 確定的影響と確率的影響の違い、線量限度の意味、低い線量のリスクについての2つの見方を説明できる。 & \pageref{ph:原-14:sec11} & \bkcb & \bkcb \\
{\sffamily\bfseries\color{\bkTT}原-15} & ウラン235 の核分裂で、1回に約 200 MeV のエネルギーと2〜3個の中性子が出ることと、連鎖反応・原子炉のしくみを、図で説明できる。 & \pageref{ph:原-15:sec1} & \bkcb & \bkcb \\
 & 増倍率 $k$・反応度・遅発中性子を使って、原子炉の出力の増え方（炉の周期）を見積もれる。 & \pageref{ph:原-15:sec6} & \bkcb & \bkcb \\
 & 4因子公式と1群の拡散方程式で、増倍率と、原子炉が臨界になる大きさを求められる。 & \pageref{ph:原-15:sec7} & \bkcb & \bkcb \\
 & 1点炉の動特性の式・反応度係数・キセノンの式を、微分方程式として読める。 & \pageref{ph:原-15:sec8} & \bkcb & \bkcb \\
 & 崩壊熱と核融合のエネルギーを数で見積もり、原子炉を止めたあとも冷やし続けなければならないわけを説明できる。 & \pageref{ph:原-15:sec11} & \bkcb & \bkcb \\
{\sffamily\bfseries\color{\bkTT}原-16} & 原子の力の使いみち（発電・医療・工業・農業・研究）と、核分裂の発見から原子炉・原子爆弾・原子力発電までの歴史を説明できる。 & \pageref{ph:原-16:sec1} & \bkcb & \bkcb \\
 & 使い終わった燃料が熱と放射線を出し続けるわけを、崩壊熱と半減期で説明し、冷やす期間と、人から離しておく期間の桁を見積もれる。 & \pageref{ph:原-16:ex2} & \bkcb & \bkcb \\
 & 福島第一原子力発電所の事故の経過と、廃炉・処理水のことを、物理の量（崩壊熱・Bq・Sv）と公的な資料の数で説明できる。 & \pageref{ph:原-16:sec5} & \bkcb & \bkcb \\
 & 放射線や原子力の数を、何を・どの基準で数えた値かに気をつけて読み、物理でわかることと、社会が選ぶことを分けて考えられる。 & \pageref{ph:原-16:sec9} & \bkcb & \bkcb \\
\end{longtblr}
\endgroup


%% ---------------- つぎの巻へ ----------------
%% 第5巻はシリーズの最後の巻。見出しは前の巻と同じ「つぎの巻へ」のまま（もくじの項目・tools/check_pages.py と合わせる）にし、
%% 中身を「全5巻のふりかえり」（5巻の色の箱）と「この先の学び」（工学の専門の科目）にした。
\bkbacktitle{つぎの巻へ}
\newcommand\bknextpill[1]{\bk@pill{#1}}
{\small 原子で、光と電子の二重性から、半導体・放射線・原子炉までを調べる道具がそろいました。物理シリーズは、この巻でおしまいです。前の4巻で「原子で」と送った約束（金属の中の電子、磁石のもと、プランクの式、とびとびのエネルギー、ダイオードとトランジスタ、相対論）は、すべてこの巻で果たしました。この先は、大学・専門の本で、それぞれの工学へ進みます。\par}\vspace{3mm}
\begin{tcolorbox}[enhanced, colframe=phVol5, colback=white, boxrule=0.9pt, arc=2.5mm, left=3mm, right=3mm, top=2mm, bottom=2mm,
  title={物理シリーズ全5巻：数える量と、同じ形の式で読みとく}, fonttitle=\sffamily\bfseries, colbacktitle=phVol5, coltitle=white]
\small
\begin{itemize}[leftmargin=1.2\zw, itemsep=2pt, topsep=1pt, label={\color{phVol5}●}]
\item {\color{phVol1}\textbf{第1巻『力学』}}：質量・運動量・エネルギーのバランス式で、動く・回る・ゆがむ・流れる物を調べた。
\item {\color{phVol2}\textbf{第2巻『熱学』}}：同じバランス式に熱とエントロピーを入れて、エンジン・冷蔵庫・発電所を調べた。
\item {\color{phVol3}\textbf{第3巻『波動』}}：ゆれと波を式で書き、共振・制御・音・光を調べた。定在波の固有値は、この巻の電子の箱へ続いた。
\item {\color{phVol4}\textbf{第4巻『電磁気学』}}：電気と磁気の法則を導き、回路・機械・電波を調べた。光が電磁波であることは、この巻の光子へ続いた。
\item {\color{phVol5}\textbf{第5巻『原子』}}：光はつぶ、電子は波。量子の考えで、原子・半導体・原子核・原子炉を調べた。
\end{itemize}
\end{tcolorbox}
\phbkhead{この巻の道具は、この先でこう使う}
{\small この巻の章と、その道具を使う大学・専門の工学の科目です（各章の「工学の駅」から集めた表）。前付の「保存則と相似の地図」は、どの科目に進んでも使える背骨です。\par}\vspace{2mm}
%% 自動生成（checks/物理第5巻/巻末/mknext_p5.py）。各章の「工学の駅」（\phstation）の科目を集めたもの。第5巻はシリーズの最後の巻なので、後ろの巻の章のかわりに、この先の工学の科目を並べる
\begin{longtblr}{colspec={Q[l,m,wd=62mm] X[l,m]}, width=\linewidth, rowhead=1, row{1}={bg=bkVol!14, font=\sffamily\bfseries\footnotesize},
  hlines={0.4pt,bkLine}, rowsep=1.5pt, cells={font=\footnotesize}}
この巻の章 & この道具を使う工学の科目 \\
{\bknextpill{原-1}\ 原子の中をのぞく：電子と原子核の発見} & 電子デバイス・真空工学・分析化学・加速器工学 \\
{\bknextpill{原-2}\ 光はつぶでもある：光電効果と光子} & 光エレクトロニクス・計測工学・放射測温・量子光学 \\
{\bknextpill{原-3}\ 電子は波でもある：ド・ブロイ波と不確定性} & 量子力学・電子顕微鏡学・材料評価・半導体工学 \\
{\bknextpill{原-4}\ 原子が出す光：エネルギー準位とスペクトル} & 原子物理学・分光分析・X線分析・レーザー工学 \\
{\bknextpill{原-5}\ 波で書く電子：シュレディンガー方程式} & 量子力学・半導体デバイス・光デバイス・量子化学 \\
{\bknextpill{原-6}\ 壁をすりぬける：トンネル効果と調和振動子} & 半導体デバイス・表面科学・分子分光学・超伝導エレクトロニクス \\
{\bknextpill{原-7}\ 電子の席の決まり：原子と周期表} & 材料工学・量子化学・磁性材料・医用画像 \\
{\bknextpill{原-8}\ 固体の中の電子：エネルギーバンド} & 固体物理・材料工学・電子材料・超伝導工学 \\
{\bknextpill{原-9}\ 一方通行の電気：pn接合とLED・太陽電池} & 電子デバイス・太陽光発電・光通信・照明工学 \\
{\bknextpill{原-10}\ スイッチと増幅：トランジスタと電子回路} & 電子回路・集積回路・パワーエレクトロニクス・計測工学 \\
{\bknextpill{原-11}\ 質量はエネルギー：相対論と原子核} & 原子核工学・加速器工学・医学物理・宇宙工学 \\
{\bknextpill{原-12}\ こわれる原子核：放射性崩壊と半減期} & 放射線計測・核医学・放射化学・年代学 \\
{\bknextpill{原-13}\ 放射線が物にあたる：放射線と物質} & 放射線工学・医学物理・遮へい設計・非破壊検査 \\
{\bknextpill{原-14}\ 放射線をはかる・身を守る：計測と防護} & 放射線管理・保健物理・医学物理・非破壊検査 \\
{\bknextpill{原-15}\ 原子炉の中の中性子：核分裂と原子炉} & 原子炉物理・原子炉工学・放射線工学・核融合工学 \\
{\bknextpill{原-16}\ 原子の力とつきあう} & 原子力工学・放射線工学・保健物理学・エネルギー工学 \\
\end{longtblr}


%% ---------------- さくいん ----------------
\bkbacktitle{さくいん}
{\small 数字はページです。その語がはじめて出てくるページや、詳しく説明しているページをのせています。}\par\vspace{2mm}
\raggedcolumns\begin{multicols}{2}\small\setlength{\parskip}{0pt}
\directlua{bk.printidx()}
\end{multicols}

%% ---------------- 奥付 ----------------
%% さくいんのあとのメモのページの柱を「メモ」にする（第1〜4巻と同じ）
\clearpage\bksetrunning{}{メモ}{メモ}%
\bkpadfour
\global\bk@decofalse\bkfolio{false}
\null\vfill
\noindent\begin{minipage}{\linewidth}\small
{\sffamily\bfseries\Large\color{bkDeep}\phSeries　第5巻}\par\vspace{1mm}
{\sffamily\bfseries\large 原子：量子から半導体・原子力へ}\par\vspace{3mm}
{\color{bkVol}\hrule height 1pt}\vspace{3mm}
\begin{tabular}{@{}l@{\hspace{6mm}}l@{}}
編者 & \phEditor \\
 & 委員長　小学校長（まとめ役） \\
 & 委員　数学・電気工学・機械工学・教育学・量子力学・建築工学・\\
 & \phantom{委員　}半導体工学の大学教授、幼稚園教諭、塾講師、\\
 & \phantom{委員　}特別支援学校教諭、心理士、科学解説者、小学校教諭 \\
デザイン & ユニバーサルデザイン担当委員 \\
挿絵 & イラスト担当委員 \\
先に読む本 & 「物理を理解するための数学」全9巻、物理シリーズ第1〜4巻 \\
版 & 第1版 Ver0（作成中）　2026年11月 \\
\end{tabular}\par\vspace{3mm}
{\color{bkVol}\hrule height 1pt}\vspace{2mm}
{\footnotesize この本の問題の答えは、巻末の「こたえ」にあります。詳しい解き方は、別冊の問題集と指導書にのせる予定です。核種のデータ・物性値の出典は、巻末の「核種のデータ・物性値・定数」の表にあります。}
\end{minipage}
\makeatother

